\section{Search architecture}\label{sec:search}
\subsection{Lloyd-relaxed CVT archive, sparse variation and baseline}
Each evolving island maintains a MAP-Elites archive~\cite{mouret2015,vassiliades2018} of \gArchiveNiches{} niches in an eight-dimensional
behaviour space (\src{crates/g5/src/g5_qd.rs}{\code{CvtMapElitesArchive}}): hook drop, inner turn angle, silhouette
complexity, grip section $H/T$, finger-groove articulation, porosity fraction, surface-texture energy, and mass. Whereas G3
used raw Halton points as niche centres, G5 runs 20 iterations of deterministic Lloyd relaxation~\cite{lloyd1982,vassiliades2018}
from Halton starts (16 fixed scrambled-Sobol samples per centre, \src{crates/g5/src/g5_qd.rs}{\code{lloyd\_cvt\_centroids}}) so that
the Voronoi cells partition $[0,1]^8$ with near-uniform volume.

In the v5.0.0 cohort every island initialized exclusively from 64 Joe--Kuo Sobol points in $[0,1]^{\gGeneCount}$~\cite{joe2008}, scrambled by an
island-seeded linear matrix and digital shift~\cite{matousek1998}. Offspring are generated by selecting two parents (with
probability $0.40$ from the 16 highest-scoring elites, otherwise by three-way curiosity tournament) and applying sparse
directional variation (\src{crates/g5/src/g5_genome.rs}{\code{reproduce\_iso\_line\_dd}}): module crossover with probability
$0.15$, self-adaptive per-module step sizes, and a sparse mask that perturbs an average of $\approx 15$ of the \gGeneCount{}
genes per child so that localized structural or geometric improvements are not scrambled by simultaneous full-genome moves.
v5.0.0 also ran random-search islands that drew every evaluation from the scrambled Sobol sequence with the same budget and
tier cascade; v\gRecord{} replaces them with the fresh-start control of Section~\ref{sec:continuation}.

\subsection{Rank-based tier promotion and bit-exact archive verification}
Fixed score gates for tier promotion either starve Tier~1 early in a run or flood it once scores rise. G5 replaces fixed
thresholds with a rolling quantile promoter (\code{RankPromotion} in \src{crates/g5/src/g5_qd.rs}{\code{g5_qd.rs}}): a window
of the last \gPromoWindow{} Tier-0 scores (every child is promoted during the first \gPromoWarmup) promotes a child to Tier~1 ($h = \gH$~mm voxel SDF, width-aware drainage flood fill, whole-blade
\code{SandwichPlate} FE, and \code{T6} ring sweep) if and only if its Tier-0 score ranks in the top $\gLaunchTierOneTopPct\,\%$ of the
rolling window.

During the search cohort, Tier~1 is the highest active tier so that compute is concentrated on whole-blade plate mechanics and
3D voxel drainage across thousands of candidates. After the cohort finishes, all \gCohortElites{} final elites across the
islands are re-evaluated at Tier~1 ($h = \gH$~mm) to verify bit-exact archive reproduction, and the top \gShortlist{} designs
are ranked at Tier~2 on the fine print grid ($h = \gHRank$~mm watertight marching-tetrahedra mesh census plus the three-case
$0.08~\mathrm{mm^2}$ \code{TET10} solid-tip submodel). Every persisted archive row stores its full \code{G5Loads} and
\code{AmConnectivityStats} structures so that \code{load\_and\_verify\_persisted\_dir} re-evaluates \code{W\_ERGOLIGHT\_v2}
bit-for-bit against the recorded score.

\subsection{Continuation, migration and one platform of record (v\gRecord)}\label{sec:continuation}
v\gRecord{} continues the v5.0.0 cohort with the decoder, requirement set and objective unchanged (decoder hash \gDecoder). Its
search runs on one platform of record: a single statically linked \code{x86\_64} binary per tool, run natively on the Linux hosts
and in a \code{linux/amd64} container under Rosetta on the macOS host. The same island gave byte-identical archives and histories
on all four hosts before any run of record, which removes the v5.0.0 limit that its \code{arm64} records were close to, but not
bit-identical with, \code{x86\_64} re-evaluations.

\emph{Warm start.} The \gVFiveElites{} final elites of v5.0.0 (its sixteen evolving islands) were re-evaluated at Tier~1 on the
platform of record and merged to the best record per niche (\src{crates/g5/src/bin/g5_continue.rs}{\code{g5-continue}}
\code{reeval} and \code{merge}): \gPoolZeroNiches{} niches, the starting pool.

\emph{Rounds and migration.} Each island of a round starts from the latest pool (its records inserted after a bit-exact
re-score) and runs \gEvalsIsland{} evaluations with its own seed. When every island of a host's round has finished, a hub merges
their archives into the latest pool, keeping the best record per niche in a recorded order, and the host's next round starts from
the new pool. Hosts therefore advance at their own speed; every pool records its sources, so the whole lineage is a deterministic
function of the recorded inputs, and the receipt recomputes it byte for byte (Section~\ref{sec:repro}).

\emph{Control.} \gControlIslands{} fresh-start islands ran \gControlEvals{} evaluations each from scrambled Sobol samples, with the
same decoder, objective and promotion and no pool or migration. They are never merged and cannot win.

\emph{Sizing on the hosts.} The continuation ran one island per core complex (four cores sharing an L3 cache) on the
Threadripper hosts, because an early pilot's time per evaluation rose steeply with the number of concurrent evaluations. A
controlled measurement after the search does not support that choice. The same \gSizingGenomes{} Tier-1 evaluations took
\gSizingOneS~s each with one evaluation per core complex and \gSizingFourS~s with four, so one per core would have given
\gSizingGain$\times$ the throughput of the setting used (\code{fleet/sizing.json}); the pilot had mixed tiers, simultaneous
multithreading and other users' load in it. The search is deterministic either way, so this cost evaluations, not correctness.
Freed memory is kept for reuse (glibc tunables: \gSizingFaultRatio$\times$ fewer page faults, identical outputs).

\emph{GPUs.} A profile of the Tier-1 evaluation puts \gSpikeEdtPct\,\% of its time in the voxel distance transform, \gSpikeConnPct\,\%
in the drainage connectivity flood fill and \gSpikePlatePct\,\% in the plate's skyline solve; Tier~1 is \gSpikeTierOneShare\,\% of an
island's time (\gSpikeTierOneS~s against \gSpikeTierZeroS~s for Tier~0, over \gSpikeChildren{} children of the starting pool). The
records of the search must stay exact CPU arithmetic, so a GPU could only shorten or skip Tier-1 work. Off-loading the voxel kernels
for free bounds Tier~1 at \gSpikeKernelBound$\times$ faster, and a perfect, free pre-screen that promoted only the true top 10\,\% of
children at \gSpikeScreenBound$\times$ more evaluations per hour, and the Tier-0 screen already keeps \gSpikeKept{} of the
\gSpikeTopTen{} best children (Spearman \gSpikeSpearman). A real predictor costs time and misses designs, so the $3\times$ bar set
for adopting a GPU path is out of reach; the search stays on the CPUs and the GPUs render the figures.
